\subsubsection{共动破盲与 Fourier--Laplace 层析}\label{subsubsec:typed-address-biaxial-completion-comoving-tomography-blindspot-resolution}
端点热核证书从 Toeplitz 矩中抽取端点原子质量；另一条破盲路径是改变读出坐标。静态 Cayley 图表把高高度偏移压到端点边界层，使径向深度呈现二次红移。共动图表随高度中心移动，把同一偏移读成一阶可见量，并允许把共动扫描剖面进一步编译成 Fourier--Laplace 层析与低秩 Hankel--Prony 证书。

\paragraph{静态二次压缩。}
令
\[
\mathbb H_R:=\{z=x+it:x>0\}
\]
为右半平面。边界 \(x=0\) 对应可见切片，偏离量为 \(\delta=x>0\)。静态 Cayley 图表为
\[
\mathcal C_0(z):=\frac{z-1}{z+1}.
\]
若
\[
z_\star=\delta+i\gamma,
\qquad
\delta>0,
\]
则
\[
a_0:=\mathcal C_0(z_\star)
=
\frac{\delta+i\gamma-1}{\delta+i\gamma+1}
\]
满足
\[
|a_0|^2
=
\frac{(\delta-1)^2+\gamma^2}{(\delta+1)^2+\gamma^2}.
\]
因此
\[
\boxed{
1-|a_0|^2
=
\frac{4\delta}{(\delta+1)^2+\gamma^2}.
}
\]
当 \(\gamma\gg1\) 且 \(\delta=O(1)\) 时，
\[
1-|a_0|^2
=
\frac{4\delta}{\gamma^2}
\bigl(1+O(\gamma^{-2})\bigr),
\qquad
1-|a_0|\asymp\frac{\delta}{\gamma^2}.
\]

\paragraph{共动 Cayley 图表。}
对高度中心 \(\tau\in\mathbb R\)，定义
\[
\boxed{
\mathcal C_\tau(z)
:=
\frac{z-i\tau-1}{z-i\tau+1}.
}
\]
它把 \(\mathbb H_R\) 映入单位圆盘，并把边界点 \(i\tau\) 送到端点 \(-1\)。同一缺陷点 \(z_\star=\delta+i\gamma\) 在共动图表中变为
\[
a_\tau:=\mathcal C_\tau(z_\star)
=
\frac{\delta+i(\gamma-\tau)-1}{\delta+i(\gamma-\tau)+1}.
\]
令
\[
E:=\gamma-\tau.
\]
则
\[
|a_\tau|^2
=
\frac{(\delta-1)^2+E^2}{(\delta+1)^2+E^2},
\]
并且
\[
\boxed{
1-|a_\tau|^2
=
\frac{4\delta}{(\delta+1)^2+(\gamma-\tau)^2}.
}
\]
若 \(|\gamma-\tau|\le E_0=O(1)\) 且 \(\delta\ll1\)，则
\[
1-|a_\tau|^2
=
\frac{4\delta}{1+E^2+O(\delta)}.
\]
特别地，在完全对齐 \(\tau=\gamma\) 时，
\[
1-|a_\gamma|^2
=
\frac{4\delta}{(1+\delta)^2}
=
4\delta+O(\delta^2).
\]

\begin{theorem}[共动图表的显影增益]\label{thm:typed-address-biaxial-completion-comoving-chart-redshift-gain}
设 \(z_\star=\delta+i\gamma\in\mathbb H_R\)，其中 \(\delta>0\)。令
\[
a_0=\mathcal C_0(z_\star),
\qquad
a_\tau=\mathcal C_\tau(z_\star).
\]
则
\[
1-|a_0|^2
=
\frac{4\delta}{(\delta+1)^2+\gamma^2},
\qquad
1-|a_\tau|^2
=
\frac{4\delta}{(\delta+1)^2+(\gamma-\tau)^2}.
\]
若 \(\gamma\to\infty\) 且 \(|\gamma-\tau|=O(1)\)，则
\[
\frac{1-|a_\tau|^2}{1-|a_0|^2}
=
\frac{(\delta+1)^2+\gamma^2}{(\delta+1)^2+(\gamma-\tau)^2}
\asymp
\gamma^2.
\]
\end{theorem}
\begin{proof}
两个公式均由 \(\mathcal C_\tau\) 的定义直接展开。比例式由分子分母比较得到。证毕。
\end{proof}

\paragraph{共动 Jensen 缺陷。}
设 \(F\) 是右半平面中的解析对象，或在各共动图表下可拉回到圆盘中的 Jensen 可积对象。定义
\[
F_\tau(w):=F(\mathcal C_\tau^{-1}(w)).
\]
对 \(0<r<1\)，置
\[
\boxed{
J_\tau(r)
:=
\frac1{2\pi}
\int_0^{2\pi}
\log|F_\tau(re^{i\theta})|\,d\theta
-\log|F_\tau(0)|.
}
\]
若 \(F_\tau\) 的圆盘零点为 \(a_{\tau,j}\)，重数为 \(m_j\)，则 Jensen 公式给出
\[
J_\tau(r)
=
\sum_{|a_{\tau,j}|<r}
m_j\log\frac{r}{|a_{\tau,j}|}.
\]
在无零半径处，
\[
\frac{d}{d\log r}J_\tau(r)
=
N_\tau(r),
\qquad
N_\tau(r):=
\sum_{|a_{\tau,j}|<r}m_j.
\]
固定最大半径 \(\rho_{\max}\) 能看见 \(z_\star\) 的条件为
\[
|a_\tau|<\rho_{\max}.
\]
利用共动公式，充分门槛可写为
\[
\frac{4\delta}{(\delta+1)^2+(\gamma-\tau)^2}
>
1-\rho_{\max}^2.
\]
若 \(|\gamma-\tau|\le E_0\)，则约化为
\[
\delta
\gtrsim
\frac{1-\rho_{\max}^2}{4}(1+E_0^2).
\]
静态图表中的相应门槛为
\[
\delta\gtrsim
\frac{1-\rho_{\max}^2}{4}\gamma^2.
\]

\begin{definition}[共动图表证书]\label{def:typed-address-biaxial-completion-comoving-chart-cert}
称
\[
\mathsf{ComovingChartCert}_{\tau,E_0}
\]
为共动图表证书，若它包含
\[
\left(
\tau,
[E^-,E^+],
E_0,
\mathcal C_\tau,
\mathsf{RationalMapLedger},
\mathsf{EndpointRef},
\mathsf{RoundHash}
\right).
\]
验证器复算
\[
\mathcal C_\tau(z)
=
\frac{z-i\tau-1}{z-i\tau+1}
\]
以及
\[
1-|\mathcal C_\tau(\delta+i\gamma)|^2
=
\frac{4\delta}{(\delta+1)^2+(\gamma-\tau)^2}.
\]
证书给出
\[
E=\gamma-\tau\in[E^-,E^+],
\qquad
|E|\le E_0.
\]
若对齐条件失败，则输出 \(\mathsf{ComovingAlignmentFail}\)。
\end{definition}

\paragraph{共动前缀预算。}
设高度窗为 \(\gamma\in[T,2T]\)。静态图表中
\[
|\partial_\gamma\mathcal C_0(\delta+i\gamma)|
\asymp
T^{-2},
\]
若要分离高度间隔 \(\Delta\gamma\)，需要
\[
2^{-b}\lesssim\frac{\Delta\gamma}{T^2},
\]
即
\[
b\gtrsim
2\log_2T+\log_2\frac1{\Delta\gamma}.
\]
共动图表中写 \(u=z-i\tau=\delta+i(\gamma-\tau)\)。在 \(|\gamma-\tau|\le E_0\) 的局部窗口中，
\[
|\partial_E\mathcal C_\tau(\delta+i\gamma)|=O(1),
\]
故局部分离只需
\[
b_{\mathrm{loc}}\gtrsim\log_2\frac1{\Delta\gamma}.
\]
若共动中心 \(\tau\) 在 \([T,2T]\) 上以步长 \(h\) 选取，则中心数量为 \(O(T/h)\)，外置图表地址需要
\[
b_{\mathrm{chart}}
\gtrsim
\log_2T+\log_2\frac1h.
\]
共动读出把静态二次深度门槛拆成图表选择预算与局部读出预算。

\begin{definition}[共动 Jensen 证书]\label{def:typed-address-biaxial-completion-comoving-jensen-cert}
称
\[
\mathsf{ComovingJensenCert}_{\tau,\rho}
\]
为共动 Jensen 证书，若它包含
\[
\left(
\mathsf{ComovingChartCert}_{\tau,E_0},
\rho,
[J_\tau(\rho)^-,J_\tau(\rho)^+],
\mathsf{CircleAverageLedger},
\mathsf{ZeroCountDerivativeLedger},
\mathsf{RoundHash}
\right).
\]
其核心核验对象为
\[
J_\tau(\rho)
=
\frac1{2\pi}
\int_0^{2\pi}
\log\left|F(\mathcal C_\tau^{-1}(\rho e^{i\theta}))\right|\,d\theta
-\log\left|F(\mathcal C_\tau^{-1}(0))\right|.
\]
对一组半径 \(0<\rho_1<\cdots<\rho_M<1\)，可用差分斜率包围环带计数：
\[
N_\tau([\rho_j,\rho_{j+1}])
\approx
\frac{J_\tau(\rho_{j+1})-J_\tau(\rho_j)}
{\log\rho_{j+1}-\log\rho_j}.
\]
\end{definition}

\paragraph{共动扫描与 Fourier--Laplace 指纹。}
令缺陷测度 \(\nu\) 定义在
\[
[0,\infty)\times\mathbb R
\]
上，其中 \(x=\delta\) 为离边界偏移，\(t=\gamma\) 为高度。共动扫描剖面抽象为
\[
\boxed{
S_\lambda(\tau)
=
\int e^{-\lambda x}K_\lambda(t-\tau)\,d\nu(x,t).
}
\]
其中 \(\lambda>0\) 是径向或 Laplace 深度参数，\(K_\lambda\) 是共动 Jensen 核、Poisson 核或差分核。对 \(\tau\) 作 Fourier 变换：
\[
\widehat S_\lambda(\omega)
=
\int_{\mathbb R}e^{-i\omega\tau}S_\lambda(\tau)\,d\tau.
\]
交换积分得到
\[
\widehat S_\lambda(\omega)
=
\widehat K_\lambda(\omega)
\int e^{-\lambda x}e^{-i\omega t}\,d\nu(x,t).
\]
定义 Fourier--Laplace 变换
\[
\boxed{
\mathcal H_\nu(\lambda,\omega)
:=
\int e^{-\lambda x}e^{-i\omega t}\,d\nu(x,t).
}
\]
于是
\[
\widehat S_\lambda(\omega)
=
\widehat K_\lambda(\omega)\mathcal H_\nu(\lambda,\omega).
\]
若 \(\widehat K_\lambda(\omega)\ne0\)，则
\[
\mathcal H_\nu(\lambda,\omega)
=
\frac{\widehat S_\lambda(\omega)}
{\widehat K_\lambda(\omega)}.
\]

\begin{theorem}[共动扫描的 Fourier 完备性]\label{thm:typed-address-biaxial-completion-comoving-scan-fourier-completeness}
设 \(\nu\) 为有限 Borel 测度，并令
\[
S_\lambda(\tau)
=
\int e^{-\lambda x}K_\lambda(t-\tau)\,d\nu(x,t).
\]
若在频率区域 \(\Omega\) 上
\[
\widehat K_\lambda(\omega)\ne0,
\qquad \omega\in\Omega,
\]
则 \(\mathcal H_\nu(\lambda,\omega)\) 在 \(\Omega\) 上由 \(S_\lambda\) 唯一决定。
\end{theorem}
\begin{proof}
由卷积 Fourier 公式，
\[
\widehat S_\lambda(\omega)
=
\widehat K_\lambda(\omega)\mathcal H_\nu(\lambda,\omega).
\]
在 \(\widehat K_\lambda(\omega)\ne0\) 的区域直接除以 \(\widehat K_\lambda(\omega)\)。若核在某些频率消失，该频率属于 kernel null band，只能输出类型化 \(\mathrm{NULL}\)。证毕。
\end{proof}

若知道所有 \(\lambda>0\) 与 \(\omega\in\mathbb R\) 上的 \(\mathcal H_\nu(\lambda,\omega)\)，则 \(\nu\) 被唯一决定。对固定 \(\lambda\)，\(\omega\mapsto\mathcal H_\nu(\lambda,\omega)\) 是加权测度 \(e^{-\lambda x}\nu(dx,dt)\) 在 \(t\)-方向的 Fourier 变换。让 \(\lambda\) 变化给出 \(x\)-方向的 Laplace 信息，而有限测度在半轴上的 Laplace 变换具有唯一性。

\begin{definition}[Fourier--Laplace 层析证书]\label{def:typed-address-biaxial-completion-fourier-laplace-holography-cert}
称
\[
\mathsf{FourierLaplaceHolographyCert}_{\Lambda,\Omega}
\]
为二维层析证书，若它给出区间 \([\lambda_{\min},\Lambda]\)、\([-\Omega,\Omega]\)、采样值 \(\{\mathcal H_{r,s}^{\pm}\}_{r,s}\)、kernel nonzero ledger、Fourier/Laplace tail ledger、inversion error ledger 与 round hash。采样点可取
\[
\lambda_r=r\Delta_\lambda,\qquad \omega_s=s\Delta_\omega.
\]
证书必须给出
\[
|\widehat K_{\lambda_r}(\omega_s)|\ge\kappa_K>0
\]
以及总误差
\[
\varepsilon_{\mathrm{tom}}
=
\varepsilon_{\mathrm{num}}+\varepsilon_F+\varepsilon_L+\varepsilon_K.
\]
\end{definition}

\paragraph{低秩缺陷与 Hankel--Prony 恢复。}
若
\[
\nu=\sum_{j=1}^{\kappa}a_j\delta_{(x_j,t_j)},\qquad a_j>0,
\]
则
\[
\mathcal H_\nu(\lambda,\omega)=
\sum_{j=1}^{\kappa}a_j e^{-\lambda x_j}e^{-i\omega t_j}.
\]
取
\[
\lambda_p=p\Delta_\lambda,\qquad
\omega_q=q\Delta_\omega,\qquad
X_j=e^{-\Delta_\lambda x_j},\quad
Y_j=e^{-i\Delta_\omega t_j},
\]
则样本为
\[
\boxed{
M_{p,q}=\sum_{j=1}^{\kappa}a_jX_j^pY_j^q.
}
\]
给定索引集 \(\mathcal A,\mathcal B\)，二维 Hankel 矩阵
\[
\mathcal H_{\mathcal A,\mathcal B}
:=
\bigl(M_{p+p',q+q'}\bigr)_{(p,q)\in\mathcal A,\ (p',q')\in\mathcal B}
\]
分解为
\[
\mathcal H_{\mathcal A,\mathcal B}
=
V_{\mathcal A}\operatorname{diag}(a_1,\ldots,a_\kappa)V_{\mathcal B}^{\mathsf T},
\qquad
(V_{\mathcal A})_{(p,q),j}=X_j^pY_j^q.
\]

\begin{theorem}[低秩缺陷的有限采样恢复]\label{thm:typed-address-biaxial-completion-low-rank-hankel-prony-recovery}
若节点 \((X_j,Y_j)\) 两两不同，且 \(V_{\mathcal A},V_{\mathcal B}\) 均满列秩 \(\kappa\)，则 \(\mathcal H_{\mathcal A,\mathcal B}\) 的秩为 \(\kappa\)。其核给出消失多项式；节点由共同零点恢复；节点确定后，权重由满秩 Vandermonde 系统唯一恢复。
\end{theorem}
\begin{proof}
由分解
\[
\mathcal H_{\mathcal A,\mathcal B}
=
V_{\mathcal A}DV_{\mathcal B}^{\mathsf T},
\qquad
D=\operatorname{diag}(a_1,\ldots,a_\kappa),
\]
以及两侧 Vandermonde 满列秩和 \(D\) 可逆，得秩为 \(\kappa\)。核向量对应的多项式 \(P(X,Y)\) 满足 \(P(X_j,Y_j)=0\)。采样阶数足够时，这些多项式的共同零点即节点集合；随后权重由线性系统唯一确定。证毕。
\end{proof}

反解公式为
\[
x_j=-\frac{\log|X_j|}{\Delta_\lambda},
\qquad
t_j\equiv-\frac{\arg Y_j}{\Delta_\omega}
\pmod{2\pi/\Delta_\omega}.
\]
避免高度混叠需要 Nyquist 条件或同余层析证书。

\begin{definition}[Hankel--Prony 证书]\label{def:typed-address-biaxial-completion-hankel-prony-cert}
称 \(\mathsf{HankelPronyCert}_{\kappa}\) 为低秩恢复证书，若它包含样本区间 \(\{M_{p,q}^{\pm}\}\)、Hankel 矩阵、interval rank cert、annihilator basis、节点盒、权重盒、Vandermonde condition ledger、residual ledger 与 round hash。验证器检查
\[
\sigma_\kappa^->0,\qquad \sigma_{\kappa+1}^+<\varepsilon_{\rm rank},
\]
并核验消失多项式、节点盒、权重盒与重构残差。配对不唯一输出 \(\mathsf{PronyPairingNull}\)，条件数过大输出 \(\mathsf{PronyConditioningNull}\)。
\end{definition}

\paragraph{有限秩 RKHS 载体。}
定义
\[
\phi_j(\lambda,\omega):=e^{-\lambda x_j}e^{-i\omega t_j}.
\]
则
\[
\mathcal H_\nu(\lambda,\omega)
=
\sum_{j=1}^{\kappa}a_j\phi_j(\lambda,\omega).
\]
核
\[
\mathcal K((\lambda,\omega),(\lambda',\omega'))
:=
\sum_{j=1}^{\kappa}
a_j e^{-(\lambda+\lambda')x_j}e^{-i(\omega-\omega')t_j}
\]
正定且秩为 \(\kappa\)，对应 RKHS 为
\[
\mathcal H_\nu^{\rm RKHS}
=
\operatorname{span}\{\phi_1,\ldots,\phi_\kappa\}.
\]
证书 \(\mathsf{RKHSRankCert}_{\kappa}\) 包含 sample Gram 矩阵、Gram PSD 区间证书、rank interval 证书、采样点、evaluation ledger 与 residual ledger，并检查
\[
G_{ab}=\mathcal K(q_a,q_b)\succeq0,\qquad
\sigma_\kappa^->0,\quad
\sigma_{\kappa+1}^+<\varepsilon.
\]

\paragraph{共动预算与融合验证。}
静态门槛
\[
\Delta_\rho\lesssim\frac{\delta}{\gamma^2}
\]
被共动链改写为
\[
\boxed{
\Delta_\rho\lesssim\frac{\delta}{1+E^2}.
}
\]
预算轴为
\[
(\Delta_\rho,\ E_\tau,\ \Omega,\ \Lambda,\ \kappa,\ L,\ \operatorname{cond}).
\]
证书 \(\mathsf{ComovingBlindBudgetCert}\) 核验半径显影、图表对齐、前缀地址、Fourier 带宽、Laplace 深度、Hankel 采样深度与 Vandermonde 条件数；任一预算不足时输出类型化 \(\mathrm{NULL}\)。

\begin{definition}[共动层析融合证书]\label{def:typed-address-biaxial-completion-comoving-tomography-fusion-cert}
共动层析融合证书
\[
\mathsf{ComovingTomographyFusionCert}
\]
由 \(\mathsf{EndpointHeatFusionCert}\)、\(\{\mathsf{ComovingChartCert}_{\tau_\ell,E_\ell}\}_\ell\)、\(\{\mathsf{ComovingJensenCert}_{\tau_\ell,\rho_j}\}_{\ell,j}\)、\(\mathsf{ScanFourierCert}\)、\(\mathsf{FourierLaplaceHolographyCert}_{\Lambda,\Omega}\)、\(\mathsf{HankelPronyCert}_{\kappa}\)、\(\mathsf{RKHSRankCert}_{\kappa}\)、\(\mathsf{ComovingBlindBudgetCert}\)、tomography reject ledger 与 transcript hash 组成。
\end{definition}

验证器 \(\mathsf{VerifyComovingTomographyFusion}\) 调用端点热核融合验证器后，核验共动图表、红移公式、对齐误差、共动 Jensen 圆周平均、扫描 Fourier 关系
\[
\widehat S_\lambda(\omega)
=
\widehat K_\lambda(\omega)\mathcal H_\nu(\lambda,\omega),
\]
核非零条件、Hankel rank、Prony 节点与权重、RKHS Gram rank 以及所有预算轴。

\begin{theorem}[共动层析融合证书的可靠性]\label{thm:typed-address-biaxial-completion-comoving-tomography-fusion-soundness}
若
\[
\mathsf{VerifyComovingTomographyFusion}
(\mathsf{ComovingTomographyFusionCert})
=
\mathrm{PASS},
\]
且理论接口已经证明坏事件 \(\mathcal B\) 强制目标区域 \(R_{\mathcal B}\) 中的缺陷质量满足
\[
\nu(R_{\mathcal B})\ge\Theta,
\]
而证书恢复出的层析测度满足
\[
\nu(R_{\mathcal B})\le\Theta^-<\Theta,
\]
则 \(\mathcal B\) 不发生。
\end{theorem}
\begin{proof}
端点热核与 Toeplitz 证书已经通过。共动图表证书保证目标高度点在某个 \(\tau_\ell\) 图表中越过半径盲区门槛。共动 Jensen 证书给出 \(S_\lambda(\tau)\) 的正确区间。Fourier--Laplace 证书由
\[
\widehat S_\lambda(\omega)
=
\widehat K_\lambda(\omega)\mathcal H_\nu(\lambda,\omega)
\]
恢复 \(\mathcal H_\nu\)，且核非零 ledger 保证除法稳定。Hankel--Prony 证书证明有限样本与 rank-\(\kappa\) 原子测度相容，并在误差预算内恢复节点与权重。因此验证器得到
\[
\nu(R_{\mathcal B})\le\Theta^-.
\]
若坏事件发生，则理论接口给出 \(\nu(R_{\mathcal B})\ge\Theta\)，矛盾。证毕。
\end{proof}

失败对象写为
\[
\mathsf{ComovingTomographyFail}
:=
(\mathrm{FailType},\Omega,\mathsf{Witness}).
\]
类型包括
\[
\mathrm{ComovingAlignmentFail},\quad
\mathrm{ComovingRadiusFail},\quad
\mathrm{KernelZeroFail},\quad
\mathrm{FourierBandwidthNull},\quad
\mathrm{LaplaceDepthNull},
\]
\[
\mathrm{HankelRankFail},\quad
\mathrm{PronyRootBoxFail},\quad
\mathrm{PronyPairingNull},\quad
\mathrm{PronyConditioningNull},\quad
\mathrm{RKHSRankFail},\quad
\mathrm{TomographyMarginFail}.
\]
端点热核路径读取边界端点原子质量 \(\mu(\{\eta\})\)；共动层析路径恢复边界层二维缺陷分布 \(\nu(dx,dt)\)。二者分别回答端点承载与边界层定位问题，并与 Jensen 半径证书组成破盲闭包。
